\documentclass[11pt, a4paper]{article}

% bfw-style.tex — gemeinsame Style-Definitionen für alle Handouts/Aufgabenhefte
% Voraussetzung: Kompilation mit xelatex (wegen fontspec)

\usepackage[ngerman]{babel}
\usepackage{geometry}
\geometry{a4paper, top=2.2cm, bottom=2.2cm, left=2.3cm, right=2.3cm}

\usepackage{fontspec}
% Fallback-Kette: Source Sans Pro -> Calibri -> Segoe UI -> Latin Modern Sans
% (Latin Modern ist immer mit MiKTeX vorhanden)
\IfFontExistsTF{Source Sans Pro}{%
  \setmainfont{Source Sans Pro}[Ligatures=TeX]%
}{\IfFontExistsTF{Calibri}{%
  \setmainfont{Calibri}[Ligatures=TeX]%
}{\IfFontExistsTF{Segoe UI}{%
  \setmainfont{Segoe UI}[Ligatures=TeX]%
}{%
  \setmainfont{Latin Modern Sans}[Ligatures=TeX]%
}}}
\IfFontExistsTF{Source Code Pro}{%
  \setmonofont{Source Code Pro}[Scale=0.92]%
}{\IfFontExistsTF{Consolas}{%
  \setmonofont{Consolas}[Scale=0.92]%
}{%
  \setmonofont{Latin Modern Mono}[Scale=0.92]%
}}

\usepackage{xcolor}
% Farbpalette: hell, freundlich, modern
\definecolor{bfwPrimary}{HTML}{0EA5A4}    % Petrol/Türkis - Primär
\definecolor{bfwPrimaryDark}{HTML}{0F766E} % dunkler Petrol für Headings
\definecolor{bfwAccent}{HTML}{F59E0B}     % Amber - Akzent
\definecolor{bfwSuccess}{HTML}{10B981}    % Grün - Erfolg / Lösung
\definecolor{bfwWarn}{HTML}{F97316}       % Orange - Achtung
\definecolor{bfwInk}{HTML}{0F172A}        % fast Schwarz
\definecolor{bfwInkSoft}{HTML}{475569}    % gedämpftes Grau
\definecolor{bfwBgSoft}{HTML}{F8FAFC}     % off-white
\definecolor{bfwBgTeal}{HTML}{ECFEFF}     % helles Türkis
\definecolor{bfwBgAmber}{HTML}{FEF3C7}    % helles Gelb
\definecolor{bfwBgRose}{HTML}{FFE4E6}     % helles Rosé für Eselsbrücken

\usepackage[colorlinks=true, linkcolor=bfwPrimaryDark, urlcolor=bfwPrimaryDark]{hyperref}
\usepackage{enumitem}
\setlist[itemize]{leftmargin=*, itemsep=2pt, topsep=2pt}
\setlist[enumerate]{leftmargin=*, itemsep=2pt, topsep=2pt}

\usepackage[most]{tcolorbox}
\usepackage{booktabs}
\usepackage{tikz}
\usetikzlibrary{decorations.pathmorphing, positioning, shapes.geometric, arrows.meta}

% Merkkästchen — wichtige Definition / Kernpunkt
\newtcolorbox{merksatz}[1][]{%
  enhanced, breakable,
  colback=bfwBgTeal,
  colframe=bfwPrimary,
  boxrule=0.6pt,
  arc=4pt,
  left=10pt, right=10pt, top=8pt, bottom=8pt,
  fonttitle=\bfseries\color{bfwPrimaryDark},
  title={\faLightbulb\ Merksatz},
  #1
}

% Eselsbrücke — Mnemoniks
\newtcolorbox{eselsbruecke}[1][]{%
  enhanced, breakable,
  colback=bfwBgRose,
  colframe=bfwAccent,
  boxrule=0.6pt,
  arc=4pt,
  left=10pt, right=10pt, top=8pt, bottom=8pt,
  fonttitle=\bfseries\color{bfwWarn},
  title={Eselsbrücke},
  #1
}

% Beispiel-Box
\newtcolorbox{beispiel}[1][]{%
  enhanced, breakable,
  colback=bfwBgSoft,
  colframe=bfwInkSoft,
  boxrule=0.4pt,
  arc=3pt,
  left=10pt, right=10pt, top=8pt, bottom=8pt,
  fonttitle=\bfseries\color{bfwInk},
  title={Beispiel},
  #1
}

% Lösungs-Box (für Aufgabenhefte)
\newtcolorbox{loesung}[1][]{%
  enhanced, breakable,
  colback=bfwBgAmber!40,
  colframe=bfwSuccess,
  boxrule=0.6pt,
  arc=3pt,
  left=10pt, right=10pt, top=8pt, bottom=8pt,
  fonttitle=\bfseries\color{bfwSuccess!70!black},
  title={Lösung},
  #1
}

% Glossar-Eintrag
\newcommand{\glossareintrag}[2]{%
  \par\noindent\textbf{\color{bfwPrimaryDark}#1} \enspace #2\par\smallskip
}

% Section-Stil — etwas farbiger als Standard
\usepackage{titlesec}
\titleformat{\section}
  {\Large\bfseries\color{bfwPrimaryDark}}
  {\thesection}{0.6em}{}
\titleformat{\subsection}
  {\large\bfseries\color{bfwPrimary}}
  {\thesubsection}{0.5em}{}
\titleformat{\subsubsection}
  {\normalsize\bfseries\color{bfwInk}}
  {\thesubsubsection}{0.4em}{}

% TikZ-Stil "skizzenhaft" — etwas weicher als Rough.js, aber stabil
\tikzset{
  roughbox/.style = {
    draw=bfwPrimaryDark, thick, fill=bfwBgTeal,
    rounded corners=4pt, inner sep=6pt
  },
  roughaccent/.style = {
    draw=bfwAccent, thick, fill=bfwBgAmber,
    rounded corners=4pt, inner sep=6pt
  },
  roughline/.style = {
    draw=bfwInkSoft, thick, -{Stealth[length=2.5mm]}
  }
}

% Bullet-Symbole (bewusst kein FontAwesome — vermeidet Font-Installations-Probleme)
\newcommand{\faLightbulb}{$\bullet$}
\newcommand{\faBookmark}{$\bullet$}

% Header/Footer
\usepackage{fancyhdr}
\pagestyle{fancy}
\fancyhf{}
\renewcommand{\headrulewidth}{0pt}
\fancyfoot[L]{\small\color{bfwInkSoft}\thepage}
\fancyfoot[R]{\small\color{bfwInkSoft} BFW M-064 Beschaffung im Büromanagement}


\title{\vspace*{-1.5cm}\Huge\bfseries\color{bfwPrimaryDark}Aufgabenheft Tag~9\\[0.4em]
  \large\color{bfwInkSoft}\textnormal{Vertiefung und Reflektion --- IHK-Mustertest themenübergreifend}}
\author{\textsf{Beschaffung im Büromanagement}\\
\textsf{\small\copyright\ Can Siebert\,\textperiodcentered\,2026}}
\date{Selbstlernmaterial \textperiodcentered{} 22.05.2026}

\begin{document}
\maketitle
\thispagestyle{empty}

\begin{merksatz}[title={\bfseries So nutzen Sie diesen Mustertest}]
Bearbeiten Sie alle Aufgaben unter Klausurbedingungen --- Stoppuhr (90~Minuten),
keine Hilfsmittel, vollständige Sätze, korrekte Operatoren. Erst danach Lösungen
ab Seite~\pageref{loesungen} vergleichen. Bei < 50\,\% Punkten gezielt im
\texttt{UEBERSICHT.md} nachlesen.
\end{merksatz}

\tableofcontents
\vspace{1em}
\hrule

\section{Aufgaben}

\subsection*{Aufgabe~1 --- Gesamtfall ,,Bürofix AG'' (12 Punkte)}

Die \emph{Bürofix~AG} bestellt am 02.\,Mai 2026 100~Bürostühle bei der
\emph{Office-Möbel KG} zum Listenpreis 200~€/Stück, Liefererrabatt 10\,\%,
Liefererskonto 2\,\%. Bezugskosten 600~€. Liefertermin: 15.\,Mai 2026.

\begin{enumerate}[label=(\alph*)]
  \item \textbf{Berechnen} Sie den \emph{Bezugspreis} pro Stuhl. \hfill (4~P.)
  \item \textbf{Erläutern} Sie, welche Schritte des Beschaffungsprozesses bis hierher durchlaufen wurden. \hfill (3~P.)
  \item \textbf{Beurteilen} Sie: Bis zum 25.\,Mai ist nichts geliefert. Welche Rechte hat die Bürofix~AG (§§ 286, 323 BGB)? \hfill (5~P.)
\end{enumerate}

\subsection*{Aufgabe~2 --- Wareneingang und Mängelrüge (8 Punkte)}

\begin{enumerate}[label=(\alph*)]
  \item \textbf{Beschreiben} Sie den Wareneingang in 4~Schritten. \hfill (4~P.)
  \item \textbf{Erläutern} Sie §\,377 HGB und seine Folge bei Versäumnis der Rügepflicht. \hfill (4~P.)
\end{enumerate}

\subsection*{Aufgabe~3 --- Käuferrechte bei Mängellieferung (10 Punkte)}

\begin{enumerate}[label=(\alph*)]
  \item \textbf{Nennen} Sie die vier Käuferrechte nach §\,437 BGB. \hfill (4~P.)
  \item \textbf{Erläutern} Sie den \emph{Vorrang der Nacherfüllung}. \hfill (3~P.)
  \item \textbf{Berechnen} Sie den geminderten Preis: Kaufpreis 600~€, Wert mangelfrei 600~€, Wert mangelhaft 480~€. \hfill (3~P.)
\end{enumerate}

\subsection*{Aufgabe~4 --- Lagerkennzahlen (10 Punkte)}

Aus der Lagerbuchhaltung der \emph{Schreibwaren-Vertrieb GmbH}: Anfangsbestand
3\,000~Stück, Endbestand 5\,000~Stück, Stückpreis 6,00~€, Jahresverbrauch
40\,000~Stück, Lagerzinssatz 8\,\%.

\begin{enumerate}[label=(\alph*)]
  \item \textbf{Berechnen} Sie den durchschnittlichen Lagerbestand und Lagerwert. \hfill (3~P.)
  \item \textbf{Berechnen} Sie LUH und Ø~Lagerdauer (360~Tage). \hfill (4~P.)
  \item \textbf{Berechnen} Sie die Lagerzinsen pro Jahr. \hfill (3~P.)
\end{enumerate}

\subsection*{Aufgabe~5 --- Eingangsrechnung und Skonto (12 Punkte)}

Sie erhalten eine Rechnung der \emph{Office-Möbel KG} über 5\,000~€ netto + 19\,\% USt.
Konditionen: ,,Zahlbar in 30~Tagen netto Kasse oder innerhalb 14~Tagen mit 2\,\% Skonto.''

\begin{enumerate}[label=(\alph*)]
  \item \textbf{Berechnen} Sie Brutto-, Skonto- und Zahlbetrag bei Skonto-Inanspruchnahme. \hfill (5~P.)
  \item \textbf{Berechnen} Sie die Effektivverzinsung der Skonto-Inanspruchnahme. \hfill (3~P.)
  \item \textbf{Beurteilen} Sie: Lohnt sich Skonto bei Bankkredit zu 9\,\%? \hfill (2~P.)
  \item \textbf{Nennen} Sie zwei Pflichtangaben einer Rechnung nach §\,14 UStG, die für den Vorsteuerabzug zwingend sind. \hfill (2~P.)
\end{enumerate}

\subsection*{Aufgabe~6 --- Eigentumsvorbehalt und Gefahrübergang (8 Punkte)}

\begin{enumerate}[label=(\alph*)]
  \item \textbf{Unterscheiden} Sie kurz die vier Formen des Eigentumsvorbehalts. \hfill (4~P.)
  \item \textbf{Erläutern} Sie §\,447 BGB und seinen Unterschied zu §\,475~II BGB. \hfill (4~P.)
\end{enumerate}

\vspace{1em}
\noindent\textbf{Punkte gesamt: 60}

\newpage

\section{Lösungen}\label{loesungen}

\subsection*{Lösung Aufgabe~1}

\begin{loesung}
\textbf{(a)} Bezugskalkulation pro Stuhl:
\begin{itemize}[itemsep=0pt]
  \item Listenpreis: 200,00~€
  \item Rabatt 10\,\%: $-20{,}00$~€ $\to$ Zieleinkaufspreis 180,00~€
  \item Skonto 2\,\%: $-3{,}60$~€ $\to$ Bareinkaufspreis 176,40~€
  \item Bezugskosten anteilig: $600 / 100 = 6{,}00$~€
  \item \textbf{Bezugspreis} $= \mathbf{182{,}40~€}$ pro Stuhl
\end{itemize}

\textbf{(b)} Schritte: Bedarfsermittlung $\to$ Bezugsquelle (Office-Möbel KG) $\to$
ggf.\ Anfrage und Angebotsvergleich $\to$ Bestellung am 02.\,Mai $\to$ wartet auf
Lieferung.

\textbf{(c)} Lieferungsverzug nach §\,286 BGB: Liefertermin (15.\,Mai) ist
\emph{kalendermäßig bestimmt} $\to$ Mahnung entbehrlich (§ 286 II Nr.\,1 BGB).
Bürofix AG kann: Erfüllung verlangen, Verzugszinsen geltend machen, eine
\emph{angemessene Nachfrist} setzen und nach Fristablauf vom Vertrag zurücktreten
(§ 323 BGB) oder Schadensersatz statt der Leistung fordern (§ 281 BGB).
\end{loesung}

\subsection*{Lösung Aufgabe~2}

\begin{loesung}
\textbf{(a)} Vier Schritte: \emph{Annahme} (Lieferung entgegennehmen) $\to$ \emph{äußere
Kontrolle} (vor dem Fahrer Anzahl/sichtbare Schäden) $\to$ \emph{innere Kontrolle}
(Auspacken, Inhalt prüfen) $\to$ \emph{Erfassung} (Buchung, Einlagerung).

\textbf{(b)} §\,377 HGB: Beim Handelskauf muss der Käufer die Ware unverzüglich
nach Ablieferung untersuchen und Mängel anzeigen. Bei Versäumnis gilt die Ware
nach §\,377 II HGB als \emph{genehmigt} --- alle Mängelrechte (Nacherfüllung,
Rücktritt, Minderung, Schadensersatz) sind verloren.
\end{loesung}

\subsection*{Lösung Aufgabe~3}

\begin{loesung}
\textbf{(a)} Vier Käuferrechte (§ 437 BGB): Nacherfüllung (§ 439 BGB), Rücktritt
(§ 323 BGB), Minderung (§ 441 BGB), Schadensersatz (§§ 280, 281 BGB).

\textbf{(b)} \emph{Vorrang der Nacherfüllung}: Der Käufer muss dem Verkäufer
zunächst Gelegenheit zur Mängelbeseitigung geben (Wahlrecht zwischen Nachbesserung
und Ersatzlieferung). Erst bei erfolglosem Ablauf einer angemessenen Frist
(meist 14~Tage) sind Rücktritt, Minderung und Schadensersatz statt der Leistung
möglich.

\textbf{(c)} Geminderter Preis $= 600 \times \frac{480}{600} = \mathbf{480{,}00~€}$.
Käufer kann \textbf{120~€} zurückfordern.
\end{loesung}

\subsection*{Lösung Aufgabe~4}

\begin{loesung}
\textbf{(a)} Ø~Lagerbestand $= (3\,000 + 5\,000) / 2 = \mathbf{4\,000~\text{Stück}}$.
Lagerwert $= 4\,000 \times 6{,}00 = \mathbf{24\,000~€}$.

\textbf{(b)} LUH $= 40\,000 / 4\,000 = \mathbf{10}$. Ø~Lagerdauer
$= 360 / 10 = \mathbf{36~\text{Tage}}$.

\textbf{(c)} Lagerzinsen $= 24\,000 \times 0{,}08 = \mathbf{1\,920~€}$ pro Jahr.
\end{loesung}

\subsection*{Lösung Aufgabe~5}

\begin{loesung}
\textbf{(a)} Bruttobetrag $= 5\,000 \times 1{,}19 = \mathbf{5\,950{,}00~€}$.
Skontobetrag (von 5\,000~€ netto) $= 5\,000 \times 0{,}02 = \mathbf{100{,}00~€}$.
Zahlbetrag $= 5\,950 - 100 = \mathbf{5\,850{,}00~€}$ (sofern Skonto nur auf netto;
Klausur prüfen!).

\textbf{(b)} Effektivverzinsung $= \frac{2 \times 360}{30 - 14} = \frac{720}{16} = \mathbf{45\,\%}$ p.\,a.

\textbf{(c)} Ja, Skonto lohnt sich: 45\,\% Effektivertrag schlagen 9\,\% Bankkredit
deutlich. Praxisregel: Skonto fast immer ziehen.

\textbf{(d)} Beispiele: Steuernummer/USt-IdNr.\ des Lieferanten, fortlaufende
Rechnungsnummer (oder: Liefer-/Leistungsdatum, Steuersatz und Steuerbetrag).
\end{loesung}

\subsection*{Lösung Aufgabe~6}

\begin{loesung}
\textbf{(a)} Vier Formen des Eigentumsvorbehalts:
\begin{itemize}[itemsep=0pt]
  \item \emph{Einfach}: Eigentum mit Bezahlung der konkreten Lieferung.
  \item \emph{Erweitert}: Eigentum erst, wenn alle Forderungen aus Geschäftsverbindung beglichen.
  \item \emph{Verlängert}: Käufer darf weiterverkaufen, Forderung wird im Voraus abgetreten.
  \item \emph{Weitergeleitet}: Vorbehalt geht auf Dritten über (z.\,B.\ Verarbeitung).
\end{itemize}

\textbf{(b)} §\,447 BGB (Versendungskauf B2B): Gefahrübergang mit Übergabe an
Spediteur/Frachtführer. \emph{§ 475~II BGB} (Verbrauchsgüterkauf): Gefahrübergang
erst mit Übergabe an den Verbraucher. Der Verbraucher trägt also nicht das
Transportrisiko, der Unternehmer schon.
\end{loesung}

\vfill
\hrule
\vspace{0.5em}
\noindent\small\color{bfwInkSoft}\copyright\ Can Siebert\,\textperiodcentered\,2026\,\textperiodcentered{} Beschaffung im Büromanagement\,\textperiodcentered{} Tag~9 Mustertest

\end{document}
